%% ----------------------------------------------------------------------------
%%  §1.3  Cenni di teoria degli insiemi
%% ----------------------------------------------------------------------------
\section{Cenni di teoria degli insiemi}
\label{sec:insiemi}

Introduciamo alcune notazioni e definizioni tratte dalla teoria degli insiemi: sono il
linguaggio con cui, da qui in avanti, scriveremo tutte le definizioni e tutti i
teoremi. Sia $S$ un insieme di natura qualsiasi. Per indicare che $x$ è un elemento di
$S$ scriveremo
\[
  x\in S \qquad (x \text{ \emph{appartiene} a } S);
\]
per indicare, invece, che $y$ non è un elemento di $S$, scriveremo
\[
  y\notin S \qquad (y \text{ \emph{non appartiene} a } S).
\]
Un insieme può essere descritto \emph{per elencazione}, scrivendone gli elementi tra
parentesi graffe, come $\{1,2,3\}$, oppure \emph{per proprietà caratteristica}, scrivendo
$\{x\in S : \mathcal{P}(x)\}$ per indicare «l'insieme degli $x$ di $S$ tali che vale la
proprietà $\mathcal{P}$». I due punti si leggono «tale che».

\begin{definizione}[Sottoinsiemi, insieme vuoto]
Se $A$ è un insieme i cui elementi sono anche elementi di $S$, diremo che $A$ è un
\emph{sottoinsieme}\index{sottoinsieme} o \emph{parte} di $S$ e scriveremo $A\subseteq S$.
In simboli,
\[
  (A\subseteq B) \ \sse\ (a\in A \implica a\in B).
\]
Tra i sottoinsiemi di $S$ si considera anche l'\emph{insieme vuoto}\index{insieme vuoto},
cioè l'insieme privo di elementi, che si indica con $\emptyset$; si conviene che
$\emptyset$ sia contenuto in ogni insieme. Se $A\subseteq B$ e $A\neq B$, si dice che
$A$ è una \emph{parte propria} di $B$ e si scrive $A\subsetneq B$. Due insiemi sono
uguali quando hanno gli stessi elementi, cioè $A=B$ se e solo se $A\subseteq B$ e
$B\subseteq A$.
\end{definizione}

\begin{osservazione}
Il simbolo $\implica$ si legge «implica», il simbolo $\sse$ si legge «se e solo se» o
«equivale». Per dimostrare che due insiemi $A$ e $B$ sono uguali si dimostra quasi
sempre la \emph{doppia inclusione}: si prende un generico $x\in A$ e si prova che
$x\in B$, poi si prende un generico $x\in B$ e si prova che $x\in A$. È una tecnica
che useremo molte volte.
\end{osservazione}

\begin{definizione}[Operazioni tra insiemi]\index{intersezione}\index{unione}\index{complementare}
Se $A$ e $B$ sono due sottoinsiemi di $S$ (\cref{fig:venn}):
\begin{itemize}
  \item l'\emph{intersezione} $A\cap B$ è l'insieme degli elementi di $S$ comuni ad $A$ e
  a $B$:
  \[ A\cap B = \{x\in S : x\in A \text{ e } x\in B\}; \]
  \item l'\emph{unione} $A\cup B$ è l'insieme degli elementi di $S$ che appartengono ad
  almeno uno dei due insiemi:
  \[ A\cup B = \{x\in S : x\in A \text{ oppure } x\in B\}; \]
  \item la \emph{differenza} (o \emph{complemento di $B$ rispetto ad $A$}) $A\setminus B$
  è l'insieme degli elementi di $A$ che non appartengono a $B$:
  \[ A\setminus B = \{x\in S : x\in A \text{ e } x\notin B\}. \]
\end{itemize}
In particolare, per $A=S$, l'insieme $S\setminus B$ si chiama \emph{complementare} di
$B$ e si indica con $B^{c}$. Evidentemente si ha
\begin{equation}\label{eq:inclusione-complementari}
  A\subseteq B \ \sse\ B^{c}\subseteq A^{c}.
\end{equation}
Due insiemi $A$ e $B$ con $A\cap B=\emptyset$ si dicono \emph{disgiunti}. L'insieme di
tutti i sottoinsiemi di $S$ si indica con $\parti{S}$ e si chiama \emph{insieme delle
parti}\index{insieme delle parti} di $S$.
\end{definizione}

\begin{figure}[htb]
\centering
\begin{tikzpicture}[scale=0.62, every node/.style={font=\footnotesize}]
  \tikzset{insA/.style={rotate=25}, insB/.style={rotate=-10}}
  \foreach \k/\titolo in {0/{$A\cap B$}, 1/{$A\cup B$}, 2/{$A\setminus B$}} {
  \begin{scope}[shift={(\k*6.4,0)}]
    \draw[grigio!60, thin] (-2.9,-2.1) rectangle (2.9,2.1);
    \node[anchor=north west, grigio] at (-2.85,2.05) {$S$};
    \begin{scope}
      \ifnum\k=0
        \clip[rotate=25] (-0.9,0) ellipse (1.8 and 1.2);
        \fill[blu!35, rotate=-10] (0.9,0) ellipse (1.8 and 1.2);
      \fi
      \ifnum\k=1
        \fill[blu!35, rotate=25] (-0.9,0) ellipse (1.8 and 1.2);
        \fill[blu!35, rotate=-10] (0.9,0) ellipse (1.8 and 1.2);
      \fi
      \ifnum\k=2
        \fill[blu!35, rotate=25] (-0.9,0) ellipse (1.8 and 1.2);
        \fill[white, rotate=-10] (0.9,0) ellipse (1.8 and 1.2);
      \fi
    \end{scope}
    \draw[blu, thick, rotate=25] (-0.9,0) ellipse (1.8 and 1.2);
    \draw[rame, thick, rotate=-10] (0.9,0) ellipse (1.8 and 1.2);
    \node[blu] at (-1.9,0.3) {$A$};
    \node[rame] at (1.9,-0.3) {$B$};
    \node at (0,-2.6) {\titolo};
  \end{scope}}
\end{tikzpicture}
\caption{Diagrammi di Venn: intersezione, unione e differenza di due sottoinsiemi $A$ e
$B$ di $S$. La regione colorata è l'insieme indicato sotto ciascun diagramma.}
\label{fig:venn}
\end{figure}

\begin{esempio}
Sia $S=\N$, $A=\{n\in\N : n \text{ è pari}\}$ e $B=\{n\in\N : n\le 7\}$. Allora
$A\cap B=\{2,4,6\}$, $A\setminus B=\{8,10,12,\dots\}$, $B\setminus A=\{1,3,5,7\}$ e
$A^{c}$ è l'insieme dei numeri dispari. L'insieme delle parti di $\{1,2\}$ è
$\parti{\{1,2\}}=\{\emptyset,\{1\},\{2\},\{1,2\}\}$: in generale, un insieme con $n$
elementi ha $2^{n}$ sottoinsiemi (lo dimostreremo nel \cref{sec:induzione}).
\end{esempio}

\begin{python}[Operazioni tra insiemi con il tipo \texttt{set}]
# Python ha un tipo nativo per gli insiemi finiti: set.
S = set(range(1, 13))               # S = {1, 2, ..., 12}
A = {n for n in S if n % 2 == 0}    # i pari, "per proprietà"
B = {n for n in S if n <= 7}
print("A ∩ B =", A & B)             # intersezione
print("A ∪ B =", A | B)             # unione
print("A \\ B =", A - B)            # differenza
print("B ⊆ A ?", B <= A)            # inclusione
print("A^c   =", S - A)             # complementare rispetto a S
\end{python}
\begin{risultato}
A ∩ B = {2, 4, 6}
A ∪ B = {1, 2, 3, 4, 5, 6, 7, 8, 10, 12}
A \ B = {8, 10, 12}
B ⊆ A ? False
A^c   = {1, 3, 5, 7, 9, 11}
\end{risultato}

\subsection{Prodotto cartesiano e relazioni}

\begin{definizione}[Prodotto cartesiano]\index{prodotto cartesiano}
Siano $A$ e $B$ due insiemi. Si chiama \emph{prodotto cartesiano} di $A$ e $B$, e si
indica con $A\times B$, l'insieme di tutte le \emph{coppie ordinate} $(a,b)$ con la
prima coordinata $a$ appartenente ad $A$ e la seconda coordinata $b$ appartenente a $B$.
Le coppie ordinate sono caratterizzate dalla seguente proprietà:
\begin{equation}\label{eq:coppie}
  (a,b)=(a',b') \quad\text{se e solo se}\quad a=a' \ \text{e}\ b=b'.
\end{equation}
\end{definizione}

\begin{marginfigure}
\centering
\begin{tikzpicture}[scale=0.55]
  \draw[-Latex, grigio] (-0.5,0) -- (4.5,0) node[below, font=\scriptsize] {$A$};
  \draw[-Latex, grigio] (0,-0.5) -- (0,3.5) node[left, font=\scriptsize] {$B$};
  \foreach \x in {1,2,3,4} \foreach \y in {1,2,3} \fill[blu] (\x,\y) circle (2pt);
  \fill[rame] (3,2) circle (3pt) node[above right, font=\scriptsize] {$(a,b)$};
  \draw[tratteggio] (3,0) node[below, font=\scriptsize] {$a$} -- (3,2) -- (0,2) node[left, font=\scriptsize] {$b$};
\end{tikzpicture}
\caption{Il prodotto cartesiano $A\times B$ di due insiemi finiti: una «griglia» di
coppie ordinate.}
\end{marginfigure}

L'aggettivo \emph{ordinate} è essenziale: $(1,2)\neq(2,1)$, mentre $\{1,2\}=\{2,1\}$.
Il prodotto cartesiano $\R\times\R$, che si indica anche con $\R^{2}$, è l'insieme delle
coppie di numeri reali: nel \cref{sec:funzioni} vedremo che esso si identifica con il
piano cartesiano, e il nome «cartesiano» viene proprio da lì.

\begin{definizione}[Relazioni di equivalenza e d'ordine]\index{relazione di equivalenza}\index{relazione d'ordine}
Nel caso particolare in cui sia $A=B$, un sottoinsieme $\mathcal{R}$ di $A\times A$ si
chiama \emph{relazione binaria} in $A$; se $(a,b)\in\mathcal{R}$ si scrive anche
$a\,\mathcal{R}\,b$ e si dice che «$a$ è in relazione con $b$». Una relazione binaria
$\mathcal{R}$ in $A$ si chiama \emph{relazione di equivalenza} se gode delle seguenti
proprietà:
\begin{enumerate}[label=\arabic*)]
  \item \emph{riflessiva}: per ogni $a\in A$ si ha $(a,a)\in\mathcal{R}$;
  \item \emph{simmetrica}: $(a,b)\in\mathcal{R}$ implica $(b,a)\in\mathcal{R}$;
  \item \emph{transitiva}: se $(a,b)\in\mathcal{R}$ e $(b,c)\in\mathcal{R}$, allora
  $(a,c)\in\mathcal{R}$.
\end{enumerate}
Se $(a,b)\in\mathcal{R}$ si scrive $a\sim b$ e si dice che $a$ e $b$ sono
\emph{equivalenti}. Si chiama invece \emph{relazione d'ordine} una relazione che gode
delle proprietà riflessiva, transitiva e
\begin{enumerate}[label=\arabic*'),start=2]
  \item \emph{antisimmetrica}: se $(a,b)\in\mathcal{R}$ e $(b,a)\in\mathcal{R}$, allora
  $a=b$.
\end{enumerate}
\end{definizione}

La relazione di minore o uguale tra numeri reali, considerata nei
\cref{sec:assiomi,sec:conseguenze}, è una relazione d'ordine: la proprietà
riflessiva segue dalla dicotomia \ref{ass:dicotomia}, quella antisimmetrica è l'assioma
\ref{ass:antisimmetrica} e quella transitiva è la \cref{prop:transitiva}. È inoltre un
ordine \emph{totale}, perché due elementi qualsiasi sono sempre confrontabili (ancora
per la dicotomia); l'inclusione $\subseteq$ tra i sottoinsiemi di un insieme è invece
un ordine \emph{parziale}: $\{1\}$ e $\{2\}$ non sono confrontabili.

\begin{definizione}[Classi di equivalenza e insieme quoziente]\index{classe di equivalenza}\index{insieme quoziente}
Data una relazione di equivalenza $\sim$ in $A$, indichiamo con $[a]$ la \emph{classe di
equivalenza} di $a\in A$, cioè l'insieme degli elementi equivalenti ad $a$:
$[a]=\{x\in A : x\sim a\}$. L'insieme delle classi di equivalenza si chiama
\emph{insieme quoziente} e si indica con $A/\!\sim$:
\[
  A/\!\sim\ = \{[a] : a\in A\}.
\]
\end{definizione}

\begin{proposizione}\label{prop:classi}
Due classi di equivalenza o coincidono o sono prive di elementi comuni. Di conseguenza
le classi di equivalenza formano una \emph{partizione} di $A$: ogni elemento di $A$
appartiene a una e una sola classe.
\end{proposizione}

\begin{schema}
\Ipotesi{$\sim$ relazione di equivalenza in $A$; $a,b\in A$.}
\Tesi{$[a]=[b]$ oppure $[a]\cap[b]=\emptyset$; ogni elemento di $A$ appartiene a una e
una sola classe.}
\end{schema}

\begin{strategia}
Dobbiamo provare un'alternativa: «o questo, o quello». La strada naturale è supporre
che \emph{non} valga la seconda possibilità, cioè che le due classi abbiano un elemento
in comune $c$, e dedurne la prima, cioè $[a]=[b]$. L'elemento $c$ fa da «ponte»: è
equivalente sia ad $a$ sia a $b$, e la transitività permette di passare da $a$ a $b$
attraverso $c$. Per l'uguaglianza di insiemi useremo la doppia inclusione.
\end{strategia}

\begin{dimostrazione}
Supponiamo che $[a]\cap[b]\neq\emptyset$ e sia $c$ un elemento comune: $c\sim a$ e
$c\sim b$. Per la proprietà simmetrica $a\sim c$; per la transitiva, da $a\sim c$ e
$c\sim b$ segue $a\sim b$. Proviamo ora che $[a]\subseteq[b]$: se $x\in[a]$, cioè
$x\sim a$, allora da $x\sim a$ e $a\sim b$ segue $x\sim b$, cioè $x\in[b]$. Scambiando
i ruoli di $a$ e $b$ (la situazione è simmetrica) si ottiene $[b]\subseteq[a]$, e
quindi $[a]=[b]$. Infine ogni $a\in A$ appartiene alla propria classe $[a]$ per la
proprietà riflessiva.
\end{dimostrazione}

\begin{esempio}[Le frazioni come classi di equivalenza]
Consideriamo l'insieme delle coppie $(m,n)$ con $m,n\in\Z$ (gli interi, che
richiameremo nel \cref{sec:nzq}) e $n\neq0$, e diciamo che
$(m,n)\sim(m',n')$ se $m\,n'=m'\,n$. Si verifica che è una relazione di equivalenza; la
classe di $(1,2)$ contiene $(2,4)$, $(3,6)$, $(-1,-2)$, \dots{}: è esattamente ciò che
intendiamo quando diciamo che $\tfrac12$, $\tfrac24$, $\tfrac36$ sono «la stessa
frazione». L'insieme quoziente è il modo rigoroso di definire i numeri razionali a
partire dagli interi; noi, avendo postulato l'esistenza di $\R$, non ne avremo bisogno,
ma è utile sapere che la costruzione è possibile.
\end{esempio}

\begin{esercizio}
Stabilire quali tra le seguenti relazioni in $\Z$ sono di equivalenza, quali d'ordine e
quali né l'una né l'altra: (a) $a\,\mathcal{R}\,b$ se $a-b$ è pari; (b)
$a\,\mathcal{R}\,b$ se $a$ divide $b$; (c) $a\,\mathcal{R}\,b$ se $a<b$.
\begin{soluzione}
(a) È di equivalenza: $a-a=0$ è pari; se $a-b$ è pari lo è anche $b-a$; se $a-b$ e
$b-c$ sono pari lo è la loro somma $a-c$. Le classi sono due: i pari e i dispari.
(b) Non è né d'ordine né di equivalenza in $\Z$: $2$ divide $-2$ e $-2$ divide $2$, ma
$2\neq-2$ (manca l'antisimmetria); e $2$ divide $4$ ma $4$ non divide $2$ (manca la
simmetria). In $\N$, invece, è una relazione d'ordine parziale. (c) Non è riflessiva
($a<a$ è falso), quindi non è né d'ordine né di equivalenza; è una relazione
d'\emph{ordine stretto}. \emph{Scelta logica:} per stabilire che una relazione
\emph{gode} di una proprietà serve una verifica valida per tutti gli elementi, mentre per
mostrare che \emph{non} ne gode basta un controesempio: per questo, nelle risposte
negative, abbiamo esibito una sola coppia.
\end{soluzione}
\end{esercizio}
